\section{はじめに}

ランサムウェア対策の多くは、エンドホスト上のプロセス、ファイルパス、
API 呼び出し、ファイル内容の差分、またはシグネチャを利用する。
これらは豊かな文脈を得られる反面、ホスト侵害、エージェント停止、
OS 差異、管理負荷の影響を受ける。これに対し、ストレージ装置は
ホストの意図を知らないが、永続データに対する read/write の時系列を
最後の防衛線として観測できる。SSD-Insider は SSD ファームウェア内で
ブロック I/O ヘッダに基づくランサムウェア検出と復旧を実証し
\cite{baek2018ssdinsider}、ストレージ内検出が現実的な研究対象であることを
示した。

本研究の焦点は、ストレージ装置内で動作可能な AutoEncoder (AE) による
異常検出である。
AE は正常な I/O 時系列だけを学習し、入力と再構成出力の乖離が大きい窓を
アノマリーとして扱う。従来のストレージレベル研究ではルール、決定木、
ランダムフォレスト、XGBoost などの教師あり分類器が有力であり、AE を
使う必然性はまだ示されていない。したがって本稿では、AE を既定解ではなく、
安価な統計と 500 KB モデルメモリ制約の下で検証されるべき仮説として扱う。

現段階の論文としての貢献は、(i) ブロック装置から見たランサムウェア
異常の分類、(ii) 10 秒統計列の固定入力契約、(iii) AE と単純規則の
公平な比較プロトコル、(iv) entropy 依存、benign workload shift、
split leakage、MNN memory fit を同時に扱う証拠ゲートである。性能値、
SOTA、実装完了、MNN 適合、500 KB 達成は、本稿の主張ではない。

本稿が扱う研究課題は以下である。

\begin{enumerate}
  \item ブロックストレージから安価に得られる 10 秒統計を
        $N$ 個並べた時系列だけで、
        ランサムウェアによる I/O 時系列変化を検出できるか。
  \item 正常学習 AE の再構成誤差は、単純なしきい値、write 比率、
        エントロピー、平均 LBA/平均転送長の変化、古典的異常検出器を上回る情報を持つか。
  \item エントロピーまたは圧縮率信号がない場合にも検出性能が残るか。
  \item MNN CPU 推論に変換したモデルが、MNN ランタイムを除く 500 KB の
        モデルメモリ内に収まるか。
  \item ファイルシステム、ボリューム利用率、暗号化済みボリューム、
        バックアップやデータベース処理などの benign workload shift に対し、
        低 false positive を維持できるか。
\end{enumerate}

これらの問いは、肯定されることだけを期待して設定するものではない。
AE が write throughput や entropy の単純規則に負ける場合、または
MNN 変換後の score ordering や 500 KB memory fit が崩れる場合、その結果も
本研究の結論として採用する。
